% ---------------------------------------------------------------------------
\section{Discussion}
\label{sec:discussion}
% ---------------------------------------------------------------------------

HyperCommon reads communities off local evidence only: a pair's shared neighbourhood,
required of all three pairs of a wedge and propagated through shared pairs, never
compared with a global null model. On ring lattices commonality does not depend on
community size (Lemma~\ref{lem:ringJ}), so one threshold suits rings of thirty and of
six hundred nodes alike, whereas on the gradient no single resolution of either Leiden
objective keeps the large rings whole and the small ones
apart~\cite{fortunato2007resolution,lancichinetti2011limits} (Table~\ref{tab:oracle}).
On the families with a few large dense rings and many small sparse ones the Constant
Potts Model at its best resolution is as accurate as HyperCommon on two worlds or
slightly more; HyperCommon's advantage there comes from its robustness at high noise
(three big, many small) and, on two worlds, from a threshold that overlapping modularity
selects better than it selects a resolution. Where HyperCommon falls short, on small
sparse rings beside large dense ones, it splits them and leaves part of their nodes
unassigned (Tables~\ref{tab:recsmall} and~\ref{tab:S-coverage-oracle}; coverage
$87\%$--$88\%$) rather than merging them with other communities: it reports only what
local evidence supports. Degree-biased noise affects it little, for which the cap of
Jaccard commonality by the degree ratio of a pair offers a mechanism: the edges that
hubs gain under preferential attachment rarely pass. Overlap needs no separate
membership rule: a shared node is kept in each of its communities through its own
admitted edges (Corollary~\ref{cor:pairdisjoint}).

Theory and benchmarks meet at the exact-recovery window: with isolated merges the
planted rings are recovered exactly at every threshold inside it. The ablations support
each test of the triple rule (Sec.~\ref{sec:ablations}): admitting open wedges, which J-thresholded
3-clique percolation does not, helps where triangles are scarce; testing the
non-adjacent pair, which the two-test variant omits, keeps overlapping communities
apart; and testing the wedge's two edges, which link communities~\cite{ahn2010link}
omit, separates HyperCommon from them most where communities overlap.

The advantage is tied to the lattice geometry rather than to unequal sizes as such
(Sec.~\ref{sec:geometry}). HyperCommon leads while communities are lattice-like
(without labels, not on the gradient) and trails once a tenth (with labels) or three
tenths (without) of each ring's edges are rewired internally, the former while the
communities are still strongly clustered; it trails on all eight LFR families. Random
geometric communities and rings with hubs also favour the Leiden objectives, while splitting the
degree of shared nodes, as LFR graphs do, leaves its lead on the unequal rings intact.
Communities joined by an admissible triple are read as one at every threshold at which
that triple and the triples it touches inside them are admitted
(Proposition~\ref{prop:sharedpairs} and the remark after it), the counterpart of the
pair-disjointness that makes overlap clean; this limits the method where overlaps are
densely connected, as in many social and collaboration
networks~\cite{yang2014structure,yang2012community}. No method is best on every
network~\cite{peel2017ground}; local triple evidence is the right description for
lattice-like communities of unequal size and density that overlap.

Five limits qualify these results. First, $\Omega$ is dominated by the largest
communities: by ONMI or by recovery averaged over all communities
(Tables~\ref{tab:onmi-oracle} and~\ref{tab:S-allcomm-oracle}), the Constant Potts Model
at its best resolution is ahead on two worlds and on three big, many small, whose small
rings it recovers better. Second, the label-free comparison applies one selector to
every method, and for partitions it is the Newman--Girvan modularity, whose resolution
limit the Constant Potts Model was introduced to avoid; it costs the Constant Potts
Model more than HyperCommon on unequal rings (from $0.71$ to $0.34$ on three big, many
small; Tables~\ref{tab:oracle} and~\ref{tab:labelfree}), so part of the label-free
margin reflects the selector. At each method's own default or selection rule,
HyperCommon still leads in 17 of the 28 ring settings, though the Constant Potts Model
at the edge density is level with it on two worlds with overlap
(Sec.~\ref{sec:planted}). Third, joining triples through shared non-adjacent pairs, which
distinguishes HyperCommon's connectivity from that of link communities, adds little: with the parameter chosen per graph the edge-connected variant is within $0.02$
in mean $\Omega$ on every ring setting and slightly more accurate on most LFR families;
we keep the non-adjacent joins
because they make every pair, not only every edge, bind at most one community
(Corollary~\ref{cor:pairdisjoint}). Fourth, on real networks HyperCommon is competitive
but rarely the most accurate (with Jaccard and its best threshold, only on college
football by F1), its mean rank on the fourteen smaller
networks is mid-field, the gains of HC$_f$ rest on choosing the function with the
labels, and the two-test variant has the higher F1 on Amazon and, with the best
threshold, on nine of the fourteen smaller networks (Sec.~\ref{sec:real}). Fifth,
HyperCommon, like every method we tested, reports many communities on clustered graphs
with no planted communities (Sec.~\ref{sec:geometry}), so a cover it returns is not by
itself evidence of community structure. Of the methods that select their own number of
groups by a statistical model, the nested stochastic block model and OSLOM were
compared: both are among the most accurate methods on LFR graphs, OSLOM the most accurate
on six of the eight families without labels, and both are behind on the ring lattices
even without overlap, where other partitions do well ($\Omega=0.50$ and $0.72$ against
$0.84$ for Leiden on uniform rings with $k=16$, and $0.27$ and $0.31$ against $0.79$ for
the Constant Potts Model on two worlds, both at $o_n=0$; Table~\ref{tab:labelfree});
BigClam was not compared (Sec.~\ref{sec:protocol}).

The threshold's best value moves with noise, so it has to be chosen per graph. Choosing
it by the overlapping modularity of the resulting cover~\cite{shen2009detect} works well
for HyperCommon: it retains at least $95\%$ of the best $\Omega$ on 27 of the 36
benchmark settings and a median $94\%$ of the best F1 on the real networks; nesting
(Theorem~\ref{thm:nesting}) makes the covers at all thresholds one hierarchy, read off a
single scoring of the wedges, so the choice costs only the evaluation of the overlapping
modularity of each cover (Sec.~\ref{sec:cost}). The criterion is global, and on the gradient it
selects too fine a cover; choosing thresholds separately on the branches of the
hierarchy, for instance by the cluster-stability criterion of
HDBSCAN~\cite{campello2013density}, is the natural refinement. Degree-ratio pruning, exact for every
threshold $t\ge0.02$, keeps the sweep moderate even with hubs: on YouTube, with
$1.1\times10^6$ nodes, it discards nine wedges in ten, and on thirty-two SNAP networks
it leaves at most $12\,m\ln m$ wedges.

The commonality function is the second degree of freedom. Within degree strata every
member of the family carries the signal, but the most accurate one varies from network
to network, sometimes by more than the choice of method changes the accuracy. On the
synthetic benchmarks Jaccard is the best single choice or close to it, and unnormalised
scores fall well behind where communities differ in degree. Learning the function and the threshold from the
structure of a given network is the natural next step.
